\begin{abstract}
\noindent
We derive the Goldfarb--Idnani dual active-set method for strictly convex
quadratic programs, prove the identities an implementation relies on, and check
each one against the running code.

The derivation is organised around a single matrix $J$ satisfying two invariants,
$J\T G J = I$ and $J\T A = [R;0]$, from which every quantity the iteration
computes follows: the primal direction is the $G$-orthogonal projection of
$G^{-1}n_*$ onto the working set's feasible subspace, and the dual direction
solves a least-squares problem in the $G^{-1}$ metric. Both admit
representation-free closed forms, so any $(J,R)$ meeting the invariants produces
the same iterate --- which licenses replacing a chain of Givens rotations by one
Householder reflection. The objective increment along a step of length $t$ is
exactly $t(t/2 + u_*)\,z\T G z$, positive in both step directions. And when the
primal cannot move and no multiplier can be reduced, the vector $r$ the solver
already holds is a Farkas certificate: the infeasibility verdict is a proof
rather than a failure to converge.

We then treat two additions to the classical method --- a primal--dual fast path
that guesses the whole active set, and warm starting across programs differing
only in the linear term --- both unguaranteed, and both sound only because strict
convexity makes the KKT conditions sufficient, so every candidate can be
certified. Neither admits a finite-termination theorem for general
$C\T x \ge b$, and the reason is structural rather than a missing anti-cycling
rule: the working-set system is a principal submatrix of $G^{-1}$ only when the
constraints are bounds. That exposure is easy to mistake for a technicality
because it never occurs in six hundred random instances, and duplicating columns
of the constraint matrix raises it to $93\%$. On real portfolio data the active
set reaches $442$ of $495$ constraints, where the synthetic families that most
of the literature tests on produce a small one.
\end{abstract}
